\begin{figure}[t]
\centering
\begin{tikzpicture}[
  node distance=0.5cm,
  stage/.style={
    draw,
    rounded corners,
    minimum height=0.7cm,
    minimum width=1.9cm,
    align=center,
    font=\scriptsize
  },
  tap/.style={
    draw,
    -Latex,
    thick,
    red!60!black
  },
  arr/.style={-Latex, thick}
]

  \node[stage] (raw)
    {Raw update\\\tiny (client-level gradient)};

  \node[stage, right=0.70cm of raw] (clip)
  {Clipped update\\\tiny (norm $\le C$)};

\node[stage, right=0.70cm of clip] (noise)
  {Clipped + noised\\\tiny update ($\sigma$ added)};

  \draw[arr] (raw) --
  node[above=-0.75cm, font=\tiny]{clip} (clip);

 \draw[arr] (clip) --
  node[above=-0.75cm, font=\tiny]{add $\mathcal{N}(0,\sigma^2)$} (noise);

  \node[tap, below=0.65cm of raw,
        font=\scriptsize, align=center] (dpfedavg)
    {\texttt{dp-fedavg}\\observes here\\\tiny \textbf{upper bound}};

  \node[tap, below=0.65cm of clip,
        font=\scriptsize, align=center] (central)
    {\texttt{central}\\observes here};

  \node[tap, below=0.65cm of noise,
        font=\scriptsize, align=center] (local)
    {\texttt{local}\\observes here\\\tiny \textbf{most protected}};

  \draw[tap] (dpfedavg) -- (raw);
  \draw[tap] (central) -- (clip);
  \draw[tap] (local) -- (noise);

\end{tikzpicture}
\caption{The three DP placements evaluated in this paper are three
points of adversary privilege along one pipeline, not three privacy
mechanisms. \texttt{dp-fedavg} sees the raw update and is the upper
bound on extractable signal; \texttt{central} sees the clipped update
before noise; \texttt{local} sees only the final, noised update. If the
upper-bound placement shows no attack signal, a null result under the
protective placements cannot be attributed to noise suppressing an
otherwise-present signal (Section~\ref{sec:framework-dp}).}
\label{fig:dp-placements}
\end{figure}